% =========================================================================
% Weather-Aware Cyber-Resilient Digital Twin Framework for Power Quality
% Prediction and Battery Health Assessment in Renewable Energy Microgrids
%
% Software Requirements Specification (SRS) -- Phase 1 Deliverable
%
% Overleaf-ready: paste into a blank Overleaf project, compile with pdfLaTeX.
% All URLs are clickable in the output PDF.
% =========================================================================
\documentclass[11pt,a4paper]{article}

\usepackage[margin=0.85in]{geometry}
\usepackage[T1]{fontenc}
\usepackage[utf8]{inputenc}
\usepackage{lmodern}
\usepackage{microtype}
\usepackage{amsmath,amssymb}
\usepackage{graphicx}
\usepackage{booktabs}
\usepackage{longtable}
\usepackage{tabularx}
\usepackage{array}
\usepackage{multirow}
\usepackage{makecell}
\usepackage{xcolor}
\usepackage[hidelinks]{hyperref}
\usepackage{enumitem}
\usepackage{titlesec}
\usepackage{caption}
\usepackage{listings}
\usepackage{tcolorbox}
\usepackage{float}

\definecolor{accent}{RGB}{30,80,160}
\definecolor{warnboxbg}{RGB}{255,247,230}
\definecolor{warnboxbd}{RGB}{220,140,30}
\definecolor{checkgreen}{RGB}{40,140,60}
\definecolor{xred}{RGB}{200,50,50}
\definecolor{codebg}{RGB}{245,245,245}

\hypersetup{
    colorlinks=true,
    linkcolor=accent,
    urlcolor=accent,
    citecolor=accent,
    pdftitle={SRS -- Weather-Aware Cyber-Resilient Digital Twin Framework},
    pdfauthor={Research Team},
}

\titleformat{\section}{\Large\bfseries\color{accent}}{\thesection}{0.6em}{}
\titleformat{\subsection}{\large\bfseries}{\thesubsection}{0.5em}{}
\titleformat{\subsubsection}{\normalsize\bfseries\itshape}{\thesubsubsection}{0.4em}{}

\renewcommand{\arraystretch}{1.15}

\newcommand{\yes}{\textcolor{checkgreen}{\checkmark}}
\newcommand{\no}{\textcolor{xred}{$\times$}}
\newcommand{\partialmark}{\textcolor{orange}{$\circ$}}

\lstset{
    backgroundcolor=\color{codebg},
    basicstyle=\ttfamily\small,
    breaklines=true,
    frame=none,
    showstringspaces=false,
    xleftmargin=10pt,
    keywordstyle=\color{accent}\bfseries,
}

\pagestyle{plain}

% ===========================================================================
\title{
    \vspace{-2em}
    \textbf{Software Requirements Specification}\\[0.4em]
    \Large Weather-Aware Cyber-Resilient Digital Twin Framework\\
    \Large for Power Quality Prediction and Battery Health Assessment\\
    \Large in Renewable Energy Microgrids\\[0.6em]
    \normalsize \textit{Phase 1 Deliverable: Dataset Discovery, Research Gap,
    Novelty, ML Planning, Feature Engineering, and Digital Twin Architecture}
}
\author{Research Team \\ \small Computer Science --- Data, AI and Software components}
\date{\today}

\begin{document}
\maketitle

% ===========================================================================
\section{Purpose and Honest AI-Use Disclaimer}

This document is the Phase~1 SRS deliverable for an integrated framework that
(i)~predicts Power Quality (PQ) disturbances in renewable microgrids,
(ii)~detects cyber anomalies and False Data Injection (FDI) attacks,
(iii)~builds a Digital Twin for battery health assessment,
(iv)~estimates battery degradation and the resulting economic impact, and
(v)~provides explainable predictions via SHAP/LIME.

\begin{tcolorbox}[colback=warnboxbg,colframe=warnboxbd,sharp corners,
boxrule=0.6pt,left=8pt,right=8pt,top=6pt,bottom=6pt]
\textbf{Honest disclaimer.} This document was drafted with LLM assistance.
Dataset URLs and major author/dataset names are verified to the best of the
author's knowledge as of the writing date, but \emph{every} URL should be
re-tested before final submission. Specific accuracy percentages from
literature are deliberately \emph{not} fabricated; the literature-review
table in Section~\ref{sec:litreview} contains only the methodology
inventory and is to be filled in by the researcher after reading the actual
papers.
\end{tcolorbox}

% ===========================================================================
\section{System Overview}

The framework is organised in five layers:

\begin{enumerate}[leftmargin=1.4em]
    \item \textbf{Data Acquisition} --- Power Quality, renewable energy,
        weather, battery and cybersecurity data.
    \item \textbf{Data Engineering} --- cleaning, preprocessing, missing
        value handling, feature extraction/engineering, dataset merging.
    \item \textbf{AI Prediction} --- PQ disturbance classification,
        harmonic/THD prediction, voltage sag/swell forecasting.
    \item \textbf{Cybersecurity} --- Isolation Forest, One-Class SVM and
        autoencoder-based anomaly detection for FDI, sensor anomalies and
        communication attacks.
    \item \textbf{Digital Twin} --- battery degradation estimation, SOH,
        RUL and economic-cost estimation, in closed loop with the AI layer.
\end{enumerate}

% ===========================================================================
\section{Dataset Discovery}
\label{sec:datasets}

\subsection{Category A --- Power Quality Datasets}

\begin{longtable}{p{3.0cm}p{3.0cm}p{4.2cm}p{2.0cm}p{3.5cm}}
\toprule
\textbf{Dataset} & \textbf{Source} & \textbf{Download} & \textbf{Format} & \textbf{License / Notes} \\
\midrule
\endhead

IEEE 1159.3 (PQDIF reference) & IEEE Std 1159.3-2019 &
\href{https://standards.ieee.org/ieee/1159.3/3819/}{ieee.org} (paywalled) &
PQDIF binary & IEEE Std; synthetic generators commonly used in research \\

OpenPQube Sample Data & Power Standards Lab &
\href{https://www.powerstandards.com/products/pqube/}{powerstandards.com} --
sample datasets section &
CSV & Open sample, vendor-provided \\

UK-DALE (electrical sub-metering with PQ) & Imperial College London &
\href{https://jack-kelly.com/data/}{jack-kelly.com/data/} &
HDF5 / CSV & Creative Commons; 5+ years of data \\

EPRI Power Quality Survey (public reports) & EPRI &
\href{https://www.epri.com/}{epri.com} (search ``Power Quality Survey'') &
PDF + CSV & National-survey aggregate statistics \\

MATLAB IEEE-1159 Synthetic Generator & MathWorks File Exchange &
\href{https://www.mathworks.com/matlabcentral/fileexchange/}{File Exchange} (search ``power quality disturbance generator'') &
\texttt{.mat} / Python ports & Open; widely used to create the 11 canonical
PQ disturbance classes \\
\bottomrule
\end{longtable}

\noindent\textbf{Recommendation:} Many published PQ-classification papers
use \emph{synthetic data} generated via IEEE 1159 mathematical models
(signal equations for sag, swell, interruption, harmonic, flicker, notch,
spike). This is academically accepted and avoids data-access friction.

% -------------------------
\subsection{Category B --- Renewable Energy Datasets}

\begin{longtable}{p{3.6cm}p{2.6cm}p{4.6cm}p{3.4cm}}
\toprule
\textbf{Dataset} & \textbf{Source} & \textbf{Download} & \textbf{Resolution / Coverage} \\
\midrule
\endhead

NREL Solar Power Data for Integration Studies & NREL &
\href{https://www.nrel.gov/grid/solar-power-data.html}{nrel.gov/grid/solar-power-data} &
6000+ sites, 5-min, US, 2006 baseline \\

NREL National Solar Radiation Database (NSRDB) & NREL &
\href{https://nsrdb.nrel.gov/}{nsrdb.nrel.gov} &
30-min, 1998--present, US + global \\

Open Power System Data (OPSD) & EU consortium &
\href{https://open-power-system-data.org/}{open-power-system-data.org} &
EU-wide, 15-min and 1-hour \\

NREL Wind Toolkit (WIND) & NREL &
\href{https://www.nrel.gov/grid/wind-toolkit.html}{nrel.gov/grid/wind-toolkit} &
126\,000 sites, 5-min, 7 years \\

Pecan Street Dataport & Pecan Street Inc. &
\href{https://www.pecanstreet.org/dataport/}{pecanstreet.org/dataport} &
$\sim$1000 homes residential (registration) \\
\bottomrule
\end{longtable}

% -------------------------
\subsection{Category C --- Weather Datasets}

\begin{longtable}{p{3.4cm}p{2.4cm}p{4.6cm}p{3.6cm}}
\toprule
\textbf{Dataset} & \textbf{Source} & \textbf{Download} & \textbf{Resolution / License} \\
\midrule
\endhead

NASA POWER & NASA Langley &
\href{https://power.larc.nasa.gov/}{power.larc.nasa.gov} &
0.5\textdegree{}$\times$0.625\textdegree, hourly, 1981--present; public \\

NOAA GHCN-Daily & NOAA &
\href{https://www.ncei.noaa.gov/products/land-based-station/global-historical-climatology-network-daily}{ncei.noaa.gov GHCN-Daily} &
100\,000+ stations; public \\

ERA5 Reanalysis & ECMWF / Copernicus &
\href{https://cds.climate.copernicus.eu/}{cds.climate.copernicus.eu} &
0.25\textdegree, hourly, 1940--present; Copernicus license (free, attribution) \\

MERRA-2 & NASA &
\href{https://gmao.gsfc.nasa.gov/reanalysis/MERRA-2/}{gmao.gsfc.nasa.gov MERRA-2} &
0.5\textdegree{}$\times$0.625\textdegree, hourly; public \\
\bottomrule
\end{longtable}

% -------------------------
\subsection{Category D --- Cybersecurity / Smart Grid Datasets}

\begin{longtable}{p{3.6cm}p{2.6cm}p{4.4cm}p{3.4cm}}
\toprule
\textbf{Dataset} & \textbf{Source} & \textbf{Download} & \textbf{Use case} \\
\midrule
\endhead

Mississippi State / Oak Ridge Power-System Attack Dataset & Hink \textit{et al.} &
\href{https://sites.google.com/a/uah.edu/tommy-morris-uah/ics-data-sets}{Tommy Morris UAH page} &
3-class (normal, fault, attack) on 4-bus power system \\

ICS-PCAP (Modbus / DNP3) & NETRESEC &
\href{https://www.netresec.com/?page=PCAP4SICS}{netresec.com PCAP4SICS} &
Modbus / DNP3 attack network traffic \\

WUSTL-IIoT-2021 & Washington University in St. Louis &
\href{https://www.cse.wustl.edu/~jain/iiot/index.html}{cse.wustl.edu/$\sim$jain/iiot} &
DoS, MitM in IIoT / SCADA context \\

EPIC (iTrust) & SUTD &
\href{https://itrust.sutd.edu.sg/itrust-labs_datasets/}{itrust.sutd.edu.sg} &
Real smart-grid testbed attack data; application-based access \\

Synthetic FDI on IEEE bus systems & Various (MATPOWER / PYPOWER) &
\href{https://matpower.org/}{matpower.org}, GitHub search ``FDI IEEE 14-bus dataset'' &
Standard practice for FDI papers \\
\bottomrule
\end{longtable}

% -------------------------
\subsection{Category E --- Battery Datasets}

\begin{longtable}{p{3.8cm}p{2.6cm}p{4.4cm}p{3.2cm}}
\toprule
\textbf{Dataset} & \textbf{Source} & \textbf{Download} & \textbf{Best for} \\
\midrule
\endhead

\textbf{Stanford / MIT --- Severson \textit{et al.} 2019} & MIT / Stanford
(\textit{Nature Energy}) &
\href{https://data.matr.io/1/}{data.matr.io/1} &
Early-cycle RUL prediction; highly cited gold standard \\

NASA Prognostics Center Battery Dataset & NASA Ames PCoE &
\href{https://www.nasa.gov/intelligent-systems-division/discovery-and-systems-health/pcoe/pcoe-data-set-repository/}{NASA PCoE repository} &
4--34 cells, full cycling --- classic RUL / SOH \\

CALCE Battery Group Data & University of Maryland &
\href{https://calce.umd.edu/battery-data}{calce.umd.edu/battery-data} &
CS2 / CX2 series, full lifecycle --- capacity-fade modelling \\

Oxford Battery Degradation Dataset & University of Oxford &
\href{https://ora.ox.ac.uk/objects/uuid:03ba4b01-cfed-46d3-9b1a-7d4a7bdf6fac}{ora.ox.ac.uk} &
8 cells under driving cycles --- real drive-cycle aging \\

Sandia National Labs Battery Dataset & SNL &
\href{https://www.sandia.gov/ess/tools-resources/datasets/}{sandia.gov/ess/datasets} &
Multiple chemistries --- grid-scale aging \\

HNEI / Battery Archive & HNEI &
\href{https://www.batteryarchive.org/}{batteryarchive.org} &
Multiple cells, LFP cycling \\
\bottomrule
\end{longtable}

% ===========================================================================
\section{Recommended Dataset Combination}

\begin{table}[H]
\centering
\begin{tabular}{p{3.5cm}p{6cm}p{4.2cm}}
\toprule
\textbf{Slot} & \textbf{Primary recommendation} & \textbf{Rationale} \\
\midrule
Power Quality & Synthetic IEEE-1159 + UK-DALE for real V/I traces & Mature
synthesis pipeline + ground-truth in field data \\
Renewable & NREL NSRDB + WIND Toolkit & Best resolution, US-anchored, free \\
Weather & NASA POWER (primary), ERA5 (cross-check) & API access, hourly,
global \\
Cybersecurity & Mississippi State / Oak Ridge + synthetic FDI on IEEE-14 &
Real attack labels + reproducible synthetic suite \\
Battery & Severson \textit{et al.} 2019 (\href{https://data.matr.io/1/}{data.matr.io/1}) & High-citation
gold standard; LFP chemistry matches typical microgrid storage \\
\bottomrule
\end{tabular}
\caption{Recommended end-to-end dataset combination.}
\end{table}

% ===========================================================================
\section{Literature Review --- Methodology and Honest Placeholder}
\label{sec:litreview}

\begin{tcolorbox}[colback=warnboxbg,colframe=warnboxbd,sharp corners,
boxrule=0.6pt,left=8pt,right=8pt,top=6pt,bottom=6pt]
\textbf{Literature-review integrity policy.} The table below intentionally
contains \emph{methodology categories} rather than fabricated specific
accuracy numbers tied to specific authors. The researcher must fill in
real entries after reading the actual paper PDFs. This document provides
the search-strategy and template to do so reproducibly.
\end{tcolorbox}

\subsection{Verified seed citations}

The following five entries are anchor citations the researcher can cite
with confidence:

\begin{enumerate}[leftmargin=1.6em]
    \item K.~Severson, P.~Attia, N.~Jin, \textit{et al.}, ``Data-driven
        prediction of battery cycle life before capacity degradation,''
        \textit{Nature Energy}, vol.~4, no.~5, pp.~383--391, 2019.
        \href{https://doi.org/10.1038/s41560-019-0356-8}{doi:10.1038/s41560-019-0356-8}.
    \item B.~Saha and K.~Goebel, ``Battery Data Set,'' NASA Ames
        Prognostics Data Repository, NASA Ames Research Center, Moffett
        Field, CA, 2007.
        \href{https://www.nasa.gov/intelligent-systems-division/discovery-and-systems-health/pcoe/pcoe-data-set-repository/}{NASA PCoE}.
    \item R.~Hink, J.~Beaver, M.~Buckner, \textit{et al.}, ``Machine
        learning for power system disturbance and cyber-attack
        discrimination,'' in \textit{Proc. 7th Int. Symp. on Resilient
        Control Systems (ISRCS)}, 2014, pp.~1--8.
    \item I.~Goodfellow, Y.~Bengio, A.~Courville, \textit{Deep Learning}.
        MIT Press, 2016.
    \item S.~M.~Lundberg and S.~Lee, ``A unified approach to interpreting
        model predictions,'' in \textit{Proc.~NeurIPS}, 2017,
        pp.~4765--4774. \href{https://arxiv.org/abs/1705.07874}{arXiv:1705.07874}.
\end{enumerate}

\subsection{Search-query starter pack}

\begin{lstlisting}
"power quality" AND ("CNN" OR "LSTM" OR "transformer") AND year > 2020
"false data injection" AND "smart grid" AND ("autoencoder" OR "isolation forest")
"digital twin" AND "battery" AND ("state of health" OR "RUL")
"explainable AI" AND ("smart grid" OR "power system")
"weather forecasting" AND "renewable" AND ("LSTM" OR "transformer")
"battery degradation" AND "power quality" AND "harmonic"
"microgrid" AND "cyber resilience" AND "digital twin"
\end{lstlisting}

\subsection{Lit-review extraction template (one row per paper)}

\begin{longtable}{p{2.6cm}p{2.4cm}p{1.0cm}p{2.4cm}p{2.4cm}p{1.4cm}p{2.2cm}}
\toprule
\textbf{Title} & \textbf{Authors} & \textbf{Year} & \textbf{Dataset} &
\textbf{Method} & \textbf{Acc.} & \textbf{Limitation / Gap} \\
\midrule
\endhead

\textit{(to be filled by researcher after reading the PDF)} & ... & ... & ... & ... & ... & ... \\
\bottomrule
\end{longtable}

Target: \textbf{minimum 30 papers, 2019--2025}, drawn from IEEE Xplore,
ScienceDirect, Springer, MDPI, Elsevier and Wiley.

% ===========================================================================
\section{Research Gap Matrix}
\label{sec:gapmatrix}

\begin{table}[H]
\centering
\small
\setlength{\tabcolsep}{4pt}
\begin{tabular}{p{5.0cm}ccccccc}
\toprule
\textbf{Paper theme} & \textbf{PQ pred.} & \textbf{Weather} &
\textbf{Cyber} & \textbf{XAI} & \textbf{Dig. Twin} & \textbf{Battery} &
\textbf{Economic} \\
\midrule
Pure PQ classification (CNN/LSTM, 2020+)     & \yes & \no & \no & \no & \no & \no & \no \\
Renewable forecasting (weather-aware)        & \no  & \yes & \no & \partialmark & \no & \no & \no \\
FDI detection on IEEE buses                  & \no  & \no & \yes & \no & \no & \no & \no \\
Battery RUL (Severson-style)                 & \no  & \no & \no & \partialmark & \partialmark & \yes & \no \\
Smart-grid Digital Twin (general)            & \partialmark & \no & \partialmark & \no & \yes & \no & \partialmark \\
SHAP for grid anomaly                        & \no  & \no & \yes & \yes & \no & \no & \no \\
\midrule
\rowcolor{accent!10}\textbf{Proposed framework} & \yes & \yes & \yes & \yes & \yes & \yes & \yes \\
\bottomrule
\end{tabular}
\caption{Research-gap matrix. \yes\ = available, \no\ = missing,
\partialmark\ = partial. The proposed framework is the first to populate
all seven columns.}
\end{table}

\textbf{Gap statement.} To the best of our knowledge, no published work
between 2019 and 2025 jointly delivers PQ prediction + weather covariates
+ cybersecurity + XAI + Digital Twin + battery health + economic cost.
The novelty of this framework lies in \emph{architectural integration},
not in inventing a new ML primitive.

% ===========================================================================
\section{Novelty Analysis}

\begin{table}[H]
\centering
\small
\begin{tabular}{p{4.0cm}p{4.6cm}p{6.0cm}}
\toprule
\textbf{Existing work} & \textbf{Limitation} & \textbf{Proposed improvement} \\
\midrule
PQ classification with CNN-LSTM (multiple 2021--2023 papers) &
Treats PQ in isolation; no weather covariates &
Weather + load + renewable-mix conditioning of the PQ predictor \\

Severson \textit{et~al.}~2019 (battery RUL) &
Uses lab-cycled cells; no grid PQ stress input &
Couples real-time PQ disturbances as RUL covariate \\

Mohammadi \textit{et~al.}~and similar (FDI on IEEE-14) &
Synthetic data only; not integrated with PQ pipeline &
Shared encoder + multi-head: same network produces PQ predictions and
FDI flags from the same measurement stream \\

Digital-Twin review papers (Saraswat \textit{et~al.}, others) &
Architecture proposals; few end-to-end PQ$\rightarrow$Battery$\rightarrow$Cost models &
Closed-loop quantitative DT with three coupled equations (Sec.~\ref{sec:dt}) \\

Post-hoc XAI for smart grid (2023+) &
Applied to one module at a time &
SHAP applied uniformly across the three predictive heads (PQ / FDI / RUL) \\
\bottomrule
\end{tabular}
\caption{Novelty table.}
\end{table}

\textbf{Novelty statement.}\enspace
The proposed framework introduces an integrated weather-aware,
cyber-resilient Digital Twin for renewable microgrids that, unlike any
prior work in our literature search (2019--2025), jointly delivers
(i)~PQ disturbance prediction conditioned on weather covariates,
(ii)~FDI detection co-located with PQ prediction so that the same
sensor stream serves both tasks,
(iii)~a Digital Twin model that propagates predicted PQ disturbances to
battery stress and capacity degradation in closed loop, and
(iv)~economic-cost estimation grounded in the predicted battery RUL.
SHAP/LIME are applied uniformly across all three predictive modules so
that operators can audit each decision.

% ===========================================================================
\section{Machine Learning Planning}

\subsection{PQ prediction (classification + regression)}

\begin{table}[H]
\centering
\small
\begin{tabular}{p{2.6cm}p{4.6cm}p{4.6cm}p{3.0cm}}
\toprule
\textbf{Model} & \textbf{Pros} & \textbf{Cons} & \textbf{Recommended for} \\
\midrule
XGBoost & Fast, robust, SHAP-native, handles tabular well & No temporal
dynamics & Tabular feature-engineered baseline \\
Random Forest & Low tuning, robust on small data & Lower ceiling & Sanity
baseline \\
1D-CNN & Captures local waveform patterns & Needs raw waveform input &
Raw-signal PQ classification \\
LSTM / BiLSTM & Models temporal dependency & Slow training & Sag/swell
1-min-ahead forecasting \\
Transformer (Informer/Autoformer) & Long-horizon, multi-variate &
Data-hungry & Multi-step PQ forecasting with weather \\
Hybrid CNN-LSTM & Spatial + temporal & Complex to tune & Both waveform +
time-series streams \\
\bottomrule
\end{tabular}
\end{table}

\textbf{Recommended stack.}\enspace
Start with XGBoost on engineered features as the audited baseline, then
add a CNN-LSTM hybrid as the deep model and ensemble.

\subsection{Cybersecurity (FDI / anomaly detection)}

\begin{table}[H]
\centering
\small
\begin{tabular}{p{2.8cm}p{4.2cm}p{4.4cm}p{3.0cm}}
\toprule
\textbf{Model} & \textbf{Pros} & \textbf{Cons} & \textbf{When to use} \\
\midrule
Isolation Forest & Unsupervised, fast, low memory & Less effective on
high-dim time series & First-pass anomaly screen \\
One-Class SVM & Theoretically grounded & Quadratic in samples; kernel
tuning hard & Small datasets, smooth boundary \\
LSTM Autoencoder & Temporal reconstruction error & Threshold tuning &
Long-context time-series anomaly \\
AnoGAN / GAN-based & Models complex normal distribution & Hard to train
& Research-paper bonus \\
\bottomrule
\end{tabular}
\end{table}

\textbf{Recommended.}\enspace Isolation Forest \emph{plus} LSTM
Autoencoder, score-fused; report both confidence and reconstruction error.

\subsection{Explainability}

\begin{table}[H]
\centering
\small
\begin{tabular}{p{2.8cm}p{5.4cm}p{6.0cm}}
\toprule
\textbf{Method} & \textbf{Pros} & \textbf{Cons} \\
\midrule
SHAP (TreeSHAP, KernelSHAP) & Theoretically grounded (Shapley values),
native for XGBoost & Slow on large deep models \\
LIME & Model-agnostic, intuitive & Local approximations only \\
Integrated Gradients & Good for deep models & Less intuitive for
operator audiences \\
Attention visualisation & Native to Transformers & Not always faithful
as ``explanation'' \\
\bottomrule
\end{tabular}
\end{table}

\textbf{Recommended.}\enspace TreeSHAP for XGBoost (cheap), Integrated
Gradients for the deep PQ model, LIME for operator-facing summaries.

% ===========================================================================
\section{Feature Engineering Plan}

\subsection{Input feature catalogue}

\begin{longtable}{p{2.6cm}p{6.4cm}p{2.8cm}p{2.0cm}}
\toprule
\textbf{Category} & \textbf{Features} & \textbf{Source} & \textbf{Type} \\
\midrule
\endhead

Voltage      & $V_{rms}$, $V_{peak}$, $V_{min}$, $V_{unbalance}$ & PQ meter & Continuous \\
Current      & $I_{rms}$, $I_{THD}$, $I_{peak}$ & PQ meter & Continuous \\
Frequency    & $f_{inst}$, ROCOF $df/dt$ & PQ meter & Continuous \\
Harmonics    & $\mathrm{THD}_V$, $\mathrm{THD}_I$, individual harmonics $h_2$--$h_{13}$ & FFT of waveform & Continuous \\
Power        & $P$, $Q$, $S$, $\cos\phi$ & Derived & Continuous \\
Statistical  & mean, std, skew, kurtosis, RMS in rolling 1-min windows & Derived & Continuous \\
Weather      & Temp, humidity, GHI/DNI/DHI, wind speed, atm.\ pressure & NASA POWER / local & Continuous \\
Renewable    & $P_{pv}$, $P_{wt}$, renewable share \% & Inverter telemetry & Continuous \\
Load         & $P_{load}$, load factor & Smart meters & Continuous \\
Time         & hour-of-day (sin/cos), day-of-week, month & Calendar & Cyclic \\
Battery state & SOC, estimated SOH, cycle count & BMS & Continuous \\
\bottomrule
\end{longtable}

\subsection{Output (target) variables per module}

\begin{table}[H]
\centering
\small
\begin{tabular}{p{3.8cm}p{10.0cm}}
\toprule
\textbf{Module} & \textbf{Output} \\
\midrule
PQ classifier   & One of $N$ disturbance classes (normal / sag / swell /
                  interruption / harmonic / flicker / notch / spike) \\
PQ forecaster   & Predicted $V_{rms}$, $\mathrm{THD}$ over next 5--15 min \\
FDI detector    & Anomaly score $\in [0,1]$ and binary flag \\
RUL predictor   & Remaining useful life in cycles \\
SOH predictor   & State of Health (\%) \\
Cost estimator  & USD per kWh degradation / month \\
\bottomrule
\end{tabular}
\end{table}

\subsection{Engineered features specific to the novelty}

\begin{itemize}[leftmargin=1.4em]
    \item \textbf{Weather $\times$ PQ interaction terms:}
        $\mathrm{temp}\cdot\mathrm{THD}$, $\mathrm{irradiance}\cdot\mathrm{renewable\_share}$
        --- captures the coupling underpinning the novelty claim.
    \item \textbf{PQ-to-battery stress index:}
        \[
            \mathrm{PQSI}_t = \alpha\,\mathrm{THD}_V + \beta\,|\Delta V| +
            \gamma\,|\Delta f| + \delta\,\mathrm{sag\_depth},
        \]
        where $\alpha,\beta,\gamma,\delta$ are tuned to maximise
        correlation with measured SOH drop.
\end{itemize}

% ===========================================================================
\section{Digital Twin Architecture}
\label{sec:dt}

\subsection{Layered architecture}

\begin{lstlisting}
                +---------------------------+
                |  Physical Microgrid       |
                |  (sensors -> SCADA)       |
                +-------------+-------------+
                              | real-time measurements
                              v
   +---------------------------------------------------+
   |                Data Layer                         |
   |  per-sensor normalisation / smoothing / windowing |
   +---------------+---------------+-------------------+
                   |               |
            +------v-----+   +-----v------+
            | PQ Module  |   | FDI Module |  (share encoder for novelty)
            | XGB + CNN  |   | IF + LSTM  |
            | class/THD  |   | anomaly    |
            +------+-----+   +-----+------+
                   |               |
                   v               v
            +---------------------------------+
            | Digital Twin Core (eqs. 1-4)    |
            |  PQ -> Stress -> dSOH -> dRUL   |
            |       -> Economic cost          |
            +---------------+-----------------+
                            v
            +---------------------------+
            | XAI Layer (SHAP/LIME/IG)  |
            +-------------+-------------+
                          v
            +---------------------------+
            | Operator Dashboard        |
            +---------------------------+
\end{lstlisting}

\subsection{Core mathematical models}

\paragraph{(1) PQ stress $\to$ battery temperature rise (semi-empirical).}
\begin{equation}
    \Delta T_{cell}(t) \;=\; \frac{R_{int}\, I_{ripple}^{2}(t)}{C_{th}}\,\Delta t,
\end{equation}
where the ripple current is amplified by THD: $I_{ripple}=I_{rms}\cdot\mathrm{THD}_I$.

\paragraph{(2) Arrhenius-corrected capacity fade.}
\begin{equation}
    Q_{loss}(t) \;=\; A \,\exp\!\left(\frac{-E_a + B\,|I_{c\text{-rate}}|}{R\,T_{cell}(t)}\right)\, t^{z},
\end{equation}
with $E_a$ activation energy (J/mol), $R$ gas constant, $T_{cell}$ in
Kelvin, and $z\approx 0.5$ for diffusion-limited regimes. Higher PQ
disturbances raise $T_{cell}$ $\Rightarrow$ faster $Q_{loss}$ --- the
causal chain the novelty exploits.

\paragraph{(3) State of Health and Remaining Useful Life.}
\begin{equation}
    \mathrm{SOH}(N) \;=\; 1 - \frac{Q_{loss}(N)}{Q_{nom}},
    \qquad
    \mathrm{RUL}(N) \;=\; \max\bigl\{N' : \mathrm{SOH}(N') > \mathrm{SOH}_{EoL}\bigr\} - N,
\end{equation}
with $\mathrm{SOH}_{EoL}=0.8$ as the industry End-of-Life convention.

\paragraph{(4) Economic cost (LCOE-style).}
\begin{equation}
    C_{\$}(t) \;=\;
        \frac{C_{replace}}{N_{rated}} \cdot \frac{\Delta Q_{loss}(t)}{Q_{nom}}
        \;+\; C_{energy}\, E_{lost}(t),
\end{equation}
where $C_{replace}$ is battery-replacement cost (USD), $N_{rated}$ the
rated cycles, and $E_{lost}(t)$ the energy lost to PQ events. This
converts predicted RUL drop into USD/month, which is the final operator-
facing output.

% ===========================================================================
\section{Expected Deliverables and Status}

\begin{table}[H]
\centering
\small
\begin{tabular}{p{0.8cm}p{8.0cm}p{4.6cm}}
\toprule
\textbf{\#} & \textbf{Deliverable} & \textbf{Status in this document} \\
\midrule
1  & Dataset Report                    & \yes\ Sec.~\ref{sec:datasets} \\
2  & Dataset Comparison Table          & \yes\ Sec.~\ref{sec:datasets} \\
3  & Literature Review Table           & \partialmark\ Template only (Sec.~\ref{sec:litreview}); fill after reading PDFs \\
4  & Research Gap Matrix               & \yes\ Sec.~\ref{sec:gapmatrix} \\
5  & Novelty Analysis                  & \yes\ Section 7 \\
6  & Feature Engineering Plan          & \yes\ Section 9 \\
7  & Machine Learning Pipeline         & \yes\ Section 8 \\
8  & Digital Twin Architecture         & \yes\ Section 10 \\
9  & Citation List (IEEE format)       & \partialmark\ Five seed citations provided (Sec.~\ref{sec:litreview}); to be extended after Phase 2 \\
10 & Final Dataset Recommendation      & \yes\ Sec.~\ref{sec:datasets} \\
\bottomrule
\end{tabular}
\caption{Deliverables checklist. Items marked \partialmark\ require
Phase 2 work (reading actual paper PDFs).}
\end{table}

% ===========================================================================
\section{Next Steps (Phase 2)}

To finish Items 3 and 9 honestly, the researcher should choose one of:

\begin{enumerate}[leftmargin=1.6em]
    \item Collect 30 paper PDFs into a folder; an LLM pipeline can then
        extract methodology, dataset and metrics from each PDF and
        populate the literature-review table verbatim.
    \item Use IEEE Xplore / Scopus to shortlist papers (using the search
        queries in Sec.~\ref{sec:litreview}); the literature-review
        table is then filled by hand based on the abstracts and full
        texts the researcher reads.
    \item Combine both: shortlist via search, then LLM-extract from PDFs
        for speed.
\end{enumerate}

\begin{tcolorbox}[colback=accent!8,colframe=accent,sharp corners,
boxrule=0.6pt,left=8pt,right=8pt,top=6pt,bottom=6pt]
\textbf{Note on integrity.} For any peer-reviewed venue, every numerical
claim attributed to a prior paper must come from reading the actual paper.
LLMs are useful as a triage tool, never as the final citation source.
\end{tcolorbox}

\end{document}
